%% ============================================================
%%  SECCIÓN 13 — Anexos
%% ============================================================

\begin{appendices}

\section{Deducción del modelo físico o relación cinemática}
\label{sec:anexo_modelo}

Desarrollo analítico y matemático del modelo físico del experimento.

\section{Propagación de incertidumbre}
\label{sec:anexo_propagacion}

A partir de la ecuación de trabajo:
\begin{equation}
    y = f(x_1, x_2, \dots, x_n)
\end{equation}

Calculando las derivadas parciales analíticas:
\begin{equation}
    \frac{\partial y}{\partial x_i} = \dots
\end{equation}

Sustituyendo en la expresión de suma en cuadratura:
\begin{equation}
    \delta y = \sqrt{\sum_{i=1}^n \left( \frac{\partial f}{\partial x_i}\,\delta x_i \right)^2}
\end{equation}

Resultados finales enmarcados:
\begin{equation}
    \boxed{y_1 = (12{,}34 \pm 0{,}56) \times 10^{-3}\ \text{unidad}}
\end{equation}

\section{Hoja de recolección de datos}
\label{sec:anexo_hoja}

\begin{figure}[H]
    \centering
    \includegraphics[width=0.80\linewidth]{uni_logo}
    \caption{Hoja de recolección de datos original firmada por el docente.}
    \label{fig:hoja_datos}
\end{figure}

\end{appendices}
